\documentclass[11pt]{article}

\usepackage[T1]{fontenc}
\usepackage{lmodern}
\usepackage[a4paper,margin=30mm]{geometry}
\usepackage{amsmath,amssymb,amsthm,mathtools}
\usepackage{booktabs}
\usepackage{microtype}
\usepackage{xcolor}
\usepackage[colorlinks=true,linkcolor=blue!45!black,
  citecolor=blue!45!black,urlcolor=blue!55!black]{hyperref}

\newtheorem{theorem}{Theorem}[section]
\newtheorem{proposition}[theorem]{Proposition}
\newtheorem{lemma}[theorem]{Lemma}
\newtheorem{corollary}[theorem]{Corollary}
\newtheorem{definition}[theorem]{Definition}
\newtheorem{remark}[theorem]{Remark}

\newcommand{\C}{\mathbb C}
\newcommand{\R}{\mathbb R}
\newcommand{\N}{\mathbb N}
\newcommand{\RH}{\mathrm{RH}}
\newcommand{\Rea}{\operatorname{Re}}
\newcommand{\Ima}{\operatorname{Im}}
\newcommand{\rad}{\operatorname{rad}}

\title{The Exact Abel Radius of the Riemann Xi Source:\
Finite Hausdorff Positivity, Dissipation, and a Narrow-Collar Selection Law}
\author{Michael Haufschild \and GPT 5.6 Sol \and Fable 5 \and Grok 4.5}
\date{}

\begin{document}
\maketitle

\begin{abstract}
We attach a positive Abel susceptibility to the completed prime--archimedean moment source of the
Riemann Xi function.  Its coefficients are the squared Euclidean norms of the canonical
Hausdorff--Bernstein packets.  The susceptibility has a unique exact open convergence radius
$R_A$, and the same source has a unique exact exponential exponent $E_*$ for both packet purity
and cumulative degree-elevation dissipation.  We prove the reciprocal law
\[
  R_AE_*=1,
\]
the unconditional bounds
\[
  \frac{250000}{251001}<R_A\le1,
  \qquad
  1\le E_*<\frac{251001}{250000},
\]
and the endpoint selection theorem
\[
  \RH\quad\Longleftrightarrow\quad R_A=1
  \quad\Longleftrightarrow\quad E_*=1.
\]
If RH fails, $E_*=B^2$ and $R_A=B^{-2}$ for an amplitude $B$ attained by an actual
off-critical Xi divisor atom, with $1<B<501/500$.

The bounds are accompanied by source-side estimate progress.  An atomwise opposite-phase
argument proves a complete $14\times44$ rectangle of calibrated Hausdorff inequalities.  Hence
all Bernstein coefficients through degree $13$ are nonnegative and the canonical negative mass
vanishes through defect degree $12$.  Independently, the Abel susceptibility converges
unconditionally on $0\le t<250000/251001$.  Thus the remaining RH-bearing estimate is the
extension of an already-convergent positive source susceptibility across the explicit collar
$250000/251001\le t<1$.  No positivity, self-adjointness, or finite spectral hypothesis is used
to assert that extension, and the paper does not claim a proof of RH.
\end{abstract}

\section{Introduction}

Connes' trace-formula programme seeks a spectral and noncommutative-geometric realization in
which the explicit formula is intrinsic and the missing positivity becomes structural
\cite{Connes1999}.  The later scaling-Hamiltonian and archimedean positivity constructions of
Connes and Consani sharpen the role of cofinal spectral and test-space limits
\cite{ConnesConsani2019,ConnesConsani2021}.  The central analytic difficulty remains unchanged:
finite positivity or finite spectral agreement does not by itself control the complete Xi
divisor.  A successful construction must transport an all-order arithmetic source and select its
critical phase in the cofinal limit.

This paper isolates one positive carrier with an exact phase boundary.  The input is the
completed prime--archimedean moment sequence, identified with the multiplicity-exact
reciprocal-square Xi divisor moments.  At every degree $n$, the sequence has a canonical real
Bernstein packet $(a_{n,i})_{0\le i\le n}$.  We use the positive packet purity
\[
  \Pi_n=\sum_{i=0}^n a_{n,i}^2
\]
and its Abel susceptibility
\[
  \chi(t)=\sum_{n\ge0}\Pi_nt^n,qquad t\ge0.
\]
The quantity $\chi$ is cancellation-resistant: the source may have signed Bernstein masses, but
each coefficient of $\chi$ is nonnegative.  It also retains the full divisor geometry.  An
off-critical zero produces an attained affine Bernstein mode of amplitude $B>1$, and that mode
fixes the exact convergence radius $B^{-2}$.

The headline result is not merely another equivalence with RH.  It contains two unconditional
estimates on the missing source quantity.  First, the completed Xi geometry confines every atom
to an ellipse of affine amplitude less than $501/500$, proving
\[
  \chi(t)<\infty\qquad
  \left(0\le t<\frac{250000}{251001}\right).
\]
Second, a calibrated opposite-phase estimate proves $616$ exact Hausdorff inequalities and hence
finite Bernstein positivity through degree $13$.  The unresolved part is localized to a collar
of exact width $1001/251001$ at the Abel boundary.

The proof uses four ingredients.  The completed source--divisor identity supplies the moments.
Hausdorff degree elevation supplies a positive Dirichlet dissipation.  A narrow sector for the
reciprocal-square divisor supplies the upper exponential envelope.  Finally, a dominant-mode
argument shows that any off-critical atom survives the complete divisor sum and supplies a
matching lower envelope.  These ingredients produce one master theorem joining finite source
positivity, exact exponential growth, exact Abel convergence, and cofinal dissipation.

\section{The completed reciprocal-square source}

Let
\[
  \xi(s)=\frac12s(s-1)\pi^{-s/2}\Gamma(s/2)\zeta(s)
\]
and let $\mathcal Z$ be its nontrivial zero multiset, with analytic multiplicity.  For each divisor
copy $\rho\in\mathcal Z$, put
\[
  a_\rho=(\rho-\tfrac12)(-i),
  \qquad
  \lambda_\rho=a_\rho^{-2}=-\frac4{(2\rho-1)^2}.
  \tag{2.1}
\]
On the critical line, $\rho=\tfrac12+i\gamma$ gives
$a_\rho=\gamma$ and $\lambda_\rho=\gamma^{-2}>0$.
The classical zero-counting estimate implies
\[
  \sum_{\rho\in\mathcal Z}|\lambda_\rho|<\infty.
  \tag{2.2}
\]
Hence every reciprocal moment below is absolutely convergent.

Let
\[
  C=2+\gamma_E-\log(4\pi),
  \qquad R_*=\frac{C}{1-C},
  \qquad A_*=\sum_{\rho\in\mathcal Z}|\lambda_\rho|.
  \tag{2.3}
\]
The completed prime--archimedean moment sequence is
\[
  \mu_n=\Rea\sum_{\rho\in\mathcal Z}\lambda_\rho^{,n+1},
  \qquad
  s_n=R_*^{-n}\mu_n.
  \tag{2.4}
\]
The normalization leaves the zeroth mass unchanged: $s_0=\mu_0>0$.

\begin{proposition}[Completed source identity]
\label{prop:source-identity}
The sequence $(\mu_n)$ is the Taylor-moment sequence of the normalized completed
prime--Gamma source germ.  For every real polynomial $p$,
\[
  \ell_\mu(p):=\sum_{n\ge0}[X^n]p\,\mu_n
  =\Rea\sum_{\rho\in\mathcal Z}\lambda_\rho p(\lambda_\rho).
  \tag{2.5}
\]
In particular, for $k,r\ge0$,
\[
  D_{k,r}:=\ell_\mu\!\left(X^k(R_*-X)^r\right)
  =\Rea\sum_{\rho\in\mathcal Z}
     \lambda_\rho^{,k+1}(R_*-\lambda_\rho)^r.
  \tag{2.6}
\]
\end{proposition}

\begin{proof}
The centered canonical product of $\xi$ has logarithmic derivative
\[
  \frac{\xi'}{\xi}\!\left(\frac12+iz\right)
  =-i\sum_{\rho\in\mathcal Z}
    \left(a_\rho^{-1}-(a_\rho-z)^{-1}\right),
\]
with locally uniform convergence away from the divisor.  Expanding at $z=0$ identifies the even
Taylor coefficients with the absolutely convergent sums
$\sum_\rho a_\rho^{-2n-2}=\sum_\rho\lambda_\rho^{n+1}$.
The completed logarithmic derivative is also the sum of the Euler logarithmic derivative and
the archimedean Gamma logarithmic derivative.  Normalizing its even Taylor coefficients gives
the prime--Gamma source germ and proves (2.4).  Linearity proves (2.5), and
$p(X)=X^k(R_*-X)^r$ gives (2.6).
\end{proof}

The endpoint geometry used below is summarized next.  It is a completed-divisor estimate, not an
assumption about the location of the zeros.

\begin{proposition}[Endpoint sector and affine-amplitude confinement]
\label{prop:sector}
For $u_\rho=\lambda_\rho/R_*$, every divisor copy satisfies
\[
  \Rea u_\rho>0,
  \qquad
  (\Ima u_\rho)^2\le\frac1{84}(\Rea u_\rho)^2.
  \tag{2.7}
\]
If $\rho$ is off the critical line, then $|u_\rho|<1/4$.  Define
\[
  \beta_\rho=|u_\rho|+|1-u_\rho|.
  \tag{2.8}
\]
Then
\[
  \beta_\rho=1\quad\text{on the critical line},
  \qquad
  1<\beta_\rho<\frac{501}{500}\quad\text{off it}.
  \tag{2.9}
\]
Moreover $\beta_\rho\to1$ along the divisor multiset.
\end{proposition}

\begin{proof}
The completed endpoint curvature identity gives $0<\Rea\lambda_\rho\le R_*$ and the quadratic
orientation estimate
$(\Ima\lambda_\rho)^2\le(\Rea\lambda_\rho)^2/84$.
The strict off-line quartet estimate gives $|\lambda_\rho|<R_*/4$.
After division by $R_*$ these are (2.7) and the stated disk bound.

If $u\in[0,1]$, then $|u|+|1-u|=1$.  Strictness in the triangle inequality gives
$|u|+|1-u|>1$ for nonreal $u$.  Combining $|u|<1/4$, $\Rea u\ge0$, and (2.7) confines $u$ to a
compact rational sector on which direct squaring gives
\[
  |u|+|1-u|<\frac{501}{500}.
\]
Finally (2.2) implies $\lambda_\rho\to0$ cofinally, hence $u_\rho\to0$ and
$\beta_\rho\to1$.
\end{proof}

\section{Bernstein packets, purity, and dissipation}

Let
\[
  B_{n,i}(X)=\binom niX^i(1-X)^{n-i},
  \qquad 0\le i\le n,
\]
and define the endpoint Bernstein packet
\[
  a_{n,i}=\ell_s(B_{n,i}),
  \qquad
  \ell_s(X^m)=s_m.
  \tag{3.1}
\]
Its total variation, purity, and negative mass are
\[
  T_n=\sum_{i=0}^n|a_{n,i}|,
  \qquad
  \Pi_n=\sum_{i=0}^na_{n,i}^2,
  \qquad
  \eta_n=\sum_{i=0}^n(a_{n,i})^-.
  \tag{3.2}
\]
Here $x^-:=\max\{-x,0\}$.  Since $\sum_iB_{n,i}=1$,
\[
  \sum_{i=0}^na_{n,i}=s_0,
  \qquad
  T_n=s_0+2\eta_n.
  \tag{3.3}
\]
Cauchy--Schwarz and the elementary $\ell^2$--$\ell^1$ comparison give
\[
  \frac{s_0^2}{n+1}\le\Pi_n\le T_n^2\le(n+1)\Pi_n.
  \tag{3.4}
\]

The source packet has the exact divisor transform
\[
  P_n(w):=\sum_{i=0}^na_{n,i}w^i
  =\sum_{\rho\in\mathcal Z}\lambda_\rho
       (1-u_\rho+u_\rho w)^n.
  \tag{3.5}
\]
The sum has real coefficients because the divisor multiset is invariant under conjugation; it is
the absolutely convergent divisor representation from Proposition~\ref{prop:source-identity}.
For $|w|=1$,
\[
  |1-u_\rho+u_\rho w|\le |1-u_\rho|+|u_\rho|=\beta_\rho.
  \tag{3.6}
\]
Consequently Proposition~\ref{prop:sector} gives
\[
  \Pi_n\le T_n^2
  \le A_*^2\left(\frac{501}{500}\right)^{2n}.
  \tag{3.7}
\]

Degree elevation is exact.  Put
\[
  \mathfrak D_n=
  \sum_{i=0}^{n}
  \frac{(n+1-i)(i+1)}{n+1}
  \bigl(a_{n+1,i}-a_{n+1,i+1}\bigr)^2.
  \tag{3.8}
\]
Then $\mathfrak D_n\ge0$.

\begin{proposition}[Exact dissipation balance]
\label{prop:dissipation}
For every $n\ge0$,
\[
  (n+2)\Pi_{n+1}=(n+1)\Pi_n+\mathfrak D_n,
  \tag{3.9}
\]
and hence
\[
  (n+1)\Pi_n=\Pi_0+\sum_{m<n}\mathfrak D_m.
  \tag{3.10}
\]
\end{proposition}

\begin{proof}
Bernstein degree elevation gives
\[
  a_{n,i}=\frac{n+1-i}{n+1}a_{n+1,i}
          +\frac{i+1}{n+1}a_{n+1,i+1}.
\]
For $\alpha+\beta=1$,
\[
  (\alpha x+\beta y)^2+\alpha\beta(x-y)^2
  =\alpha x^2+\beta y^2.
\]
Apply this identity with
$\alpha=(n+1-i)/(n+1)$ and $\beta=(i+1)/(n+1)$, multiply by $n+1$, and sum over $i$.
The column weights sum to $n+2$, giving (3.9).  Telescoping gives (3.10).
\end{proof}

\section{The exact Abel radius}

\begin{definition}
The positive endpoint susceptibility is
\[
  \chi(t)=\sum_{n=0}^\infty\Pi_nt^n,
  \qquad t\ge0.
  \tag{4.1}
\]
An \emph{exact Abel radius} is a number $R_A>0$ such that $\chi(t)<\infty$ for
$0\le t<R_A$ and $\chi(t)=\infty$ for $t>R_A$.
\end{definition}

\begin{definition}
A number $E>0$ is an \emph{exact purity exponent} if
\[
  \frac{\Pi_n}{r^n}\longrightarrow0\quad(E<r),
  \qquad
  \frac{\Pi_n}{r^n}\longrightarrow+\infty\quad(0<r<E).
  \tag{4.2}
\]
The analogous definition for the cumulative dissipation uses the left side of (3.10).
\end{definition}

\begin{lemma}[Source-attained lower envelope]
\label{lem:lower-envelope}
If RH is false, there are $B,\delta>0$ with
\[
  1<B<\frac{501}{500}
\]
such that $B=\beta_{\rho_0}$ for an actual divisor copy $\rho_0$, every
$\beta_\rho\le B$, and eventually
\[
  \delta B^{2n}\le(n+1)\Pi_n.
  \tag{4.3}
\]
For every $n$ one also has
\[
  \Pi_n\le A_*^2B^{2n}.
  \tag{4.4}
\]
\end{lemma}

\begin{proof}
An off-critical copy has $\beta_\rho>1$ by Proposition~\ref{prop:sector}, while
$\beta_\rho\to1$.  Therefore the supremum $B>1$ is attained on a finite nonempty packet, and
$B<501/500$.  Inequality (4.4) is (3.6) followed by absolute summation and
$\Pi_n\le T_n^2$.

Choose a maximal copy $\rho_0$ and the unique unit phase $w$ at which
$|1-u_{\rho_0}+u_{\rho_0}w|=B$.  Equality rigidity for the Bernstein ellipse says that any other
copy attaining modulus $B$ at this same phase has either $u_\rho=u_{\rho_0}$ or the reflected
coordinate $u_\rho=1-\overline{u_{\rho_0}}$.  The latter alternative is impossible because every
maximal copy is off the line and satisfies $|u_\rho|<1/4$.  Hence all terms attaining modulus $B$
have the same affine base and the same reciprocal-square coefficient.  Their grouped coefficient
is a positive integer multiple of the nonzero number $\lambda_{\rho_0}$, and therefore is
nonzero.  Every remaining affine base has modulus strictly below $B$.  After division by $B^n$,
the complementary divisor sum tends to zero by dominated convergence, using
$\sum_\rho|\lambda_\rho|<\infty$.  Thus, for some $c>0$,
\[
  cB^n\le |P_n(w)|
\]
eventually.  Cauchy--Schwarz gives
$|P_n(w)|^2\le(n+1)\Pi_n$, proving (4.3) with $\delta=c^2$.
\end{proof}

\begin{theorem}[Finite-to-cofinal Abel selection master theorem]
\label{thm:master}
The completed Xi source has the following unconditional properties.

\begin{enumerate}
\item The calibrated source differences satisfy
\[
  D_{k,r}\ge0
  \qquad(0\le k\le13,\ 0\le r\le43).
  \tag{4.5}
\]
Consequently $a_{n,i}\ge0$ for every $n\le13$ and every $i$, and the canonical Bernstein
negative mass vanishes through defect degree $12$.

\item There is a unique exact Abel radius $R_A$ and a unique exact purity exponent $E_*$.  The
same $E_*$ is the exact exponent of the cumulative dissipation in (3.10), and
\[
  R_AE_*=1.
  \tag{4.6}
\]

\item The exact parameters lie in the narrow chambers
\[
  \frac{250000}{251001}<R_A\le1,
  \qquad
  1\le E_*<\frac{251001}{250000}.
  \tag{4.7}
\]
In particular,
\[
  \chi(t)<\infty
  \qquad\left(0\le t<\frac{250000}{251001}\right).
  \tag{4.8}
\]

\item The critical phase is selected exactly by
\[
  \RH
  \Longleftrightarrow R_A=1
  \Longleftrightarrow E_*=1
  \Longleftrightarrow
  \chi(t)<\infty\quad(0\le t<1).
  \tag{4.9}
\]
If RH is false, then $E_*=B^2$ and $R_A=B^{-2}$ for the source-attained amplitude in
Lemma~\ref{lem:lower-envelope}.

\item The same phase is equivalently the cofinally linearly dissipative phase:
\[
  \RH\Longleftrightarrow
  \exists K>0\ \forall N\ \exists n\ge N,
  \qquad
  \Pi_0+\sum_{m<n}\mathfrak D_m\le K(n+1).
  \tag{4.10}
\]
\end{enumerate}
\end{theorem}

\begin{proof}
We first prove (4.5).  Put $z=\lambda_\rho$ and $w=R_*-z$.  On the critical line both are
nonnegative real.  Off the line they lie in the open right half-plane and have opposite imaginary
parts.  Proposition~\ref{prop:sector} gives the rational angle bounds
\[
  |\arg z|\le\frac{100}{913},
  \qquad
  |\arg w|\le\frac{100}{2739}.
  \tag{4.11}
\]
Because the two arguments have opposite signs,
\[
  |(k+1)\arg z+r\arg w|
  \le\max\{(k+1)|\arg z|,r|\arg w|\}.
\]
For $k\le13$ and $r\le43$, the right side is at most
\[
  \max\left\{\frac{1400}{913},\frac{4300}{2739}\right\}<\frac\pi2.
\]
Thus $\Rea(z^{k+1}w^r)\ge0$ atom by atom.  Absolute summation in (2.6) proves all
$14\cdot44=616$ inequalities.

For $n\le13$ and $0\le i\le n$, the Bernstein coefficient is a positive binomial multiple of
the calibrated difference with indices $(i,n-i)$.  Hence it is nonnegative.  Defect degree $N$
uses the degree-$N+1$ packet, so its negative mass is zero for $N\le12$.

Next, (3.7) and geometric comparison prove (4.8), since
\[
  \left(\frac{501}{500}\right)^2t<1
  \quad\text{when}\quad
  t<\frac{250000}{251001}.
  \tag{4.12}
\]
Under RH every $u_\rho$ is real and lies in $[0,1]$.  Thus every Bernstein mass is nonnegative,
$T_n=s_0$, and
\[
  \frac{s_0^2}{n+1}\le\Pi_n\le s_0^2.
  \tag{4.13}
\]
It follows that the exact exponent is $E_*=1$ and that $\chi(t)$ converges for every $t<1$.
The lower bound in (4.13) excludes any exact Abel radius larger than one, so $R_A=1$.

Suppose now that RH is false and take $B,\delta$ from Lemma~\ref{lem:lower-envelope}.  The upper
and lower envelopes (4.3)--(4.4) show directly that $E_*=B^2$: normalization above $B^2$ is
dominated by a strict geometric contraction, while normalization below $B^2$ diverges even after
division by $n+1$.  They also identify the Abel radius.  If $0\le t<B^{-2}$, then
\[
  \Pi_nt^n\le A_*^2(B^2t)^n
\]
is summable.  If $t>B^{-2}$, choose $r=t^{-1}<B^2$ in the lower exact-exponent law; then
$\Pi_nt^n\to+\infty$, so the terms of the Abel series do not tend to zero.  Hence
$R_A=B^{-2}$.

The midpoint argument proves uniqueness of both exact parameters.  Two distinct exact exponents
would impose zero and infinite limits at an intermediate normalization; two distinct Abel radii
would make $\chi$ finite and infinite at an intermediate parameter.  The case analysis above
gives $R_AE_*=1$ and the strict bounds (4.7).

Finally multiplication by $n+1$ does not change an exact exponential exponent.  Equation (3.10)
therefore gives the same $E_*$ for cumulative dissipation.  Under RH, (4.13) and (3.10) give a
linear bound at every degree.  If RH is false, Lemma~\ref{lem:lower-envelope} forces
$\Pi_n\to+\infty$; hence no fixed linear bound can recur cofinally in (3.10).  This proves
(4.10) and completes the theorem.
\end{proof}

\section{The positive Connes carrier obligation}

Theorem~\ref{thm:master} turns the remaining estimate into a positive carrier specification.
The completed arithmetic source already produces a positive susceptibility on a nearly critical
interval, and the divisor itself fixes its exact boundary.  A successful Connes-type realization
must therefore carry the same packet purity or an exactly identified equivalent observable and
establish the unit-boundary extension
\[
  \chi(t)<\infty\qquad(0\le t<1).
  \tag{5.1}
\]
Equivalently, it must produce the cofinal source-coupled budget (4.10).  This is weaker than a
uniform bound at every degree and stronger than finite positivity, finite self-adjointness, or
finite spectral agreement.

\begin{corollary}[Narrow-collar selection principle]
\label{cor:collar}
To prove RH through the endpoint Bernstein carrier it is sufficient and necessary to extend the
unconditional convergence in (4.8) across
\[
  \frac{250000}{251001}\le t<1.
  \tag{5.2}
\]
Any proof of this extension must use source--divisor incidence capable of excluding an attained
off-critical amplitude $1<B<501/500$.
\end{corollary}

\begin{proof}
The first statement is (4.8)--(4.9).  If RH fails, Theorem~\ref{thm:master} gives an actual
divisor amplitude $B$ and an interior divergence point above $R_A=B^{-2}$.  Thus an argument that
does not distinguish the source-attained transverse amplitude cannot establish (5.1).
\end{proof}

This formulation also separates what has and has not been estimated.  The finite rectangle proves
an exact source-positive southwest block.  The sector theorem proves an open interval of
susceptibility convergence.  The dissipation identity proves that only one recurring linear
budget is needed.  None of these facts implies the collar extension: that extension is precisely
the remaining RH-bearing theorem.

The carrier is cofinal in two senses.  It uses every Bernstein degree, and its source moments use
the full multiplicity-exact divisor.  This is the structural feature needed by a semilocal trace
construction: the finite stages may be positive, but their normalization must converge in a way
that controls the actual all-order source susceptibility.  In this form the required object is
not merely ``not finite''.  It must be a cofinal, divisor-preserving, source-coupled section whose
positive energy obeys (4.10), or equivalently whose Abel transform obeys (5.1).

\section{Limitations and next theorem}

The paper does not prove the collar extension (5.1) and therefore does not prove RH.  The exact
radius theorem is stronger than a qualitative reformulation because it proves a strict
unconditional lower bound for the radius, a finite block of exact source positivity, and a common
exponent for purity and dissipation.  Nevertheless, selecting the endpoint value $R_A=1$ remains
equivalent to RH.

The next theorem must estimate the actual completed prime--Gamma source, not only the enclosing
sector.  Any one of the following source statements would complete the route:
\begin{enumerate}
\item $\chi(t)<\infty$ for every $t<1$;
\item a cofinal linear bound for the cumulative dissipation in (4.10);
\item a cofinal bounded packet-purity subsequence;
\item a uniformly bounded aggregate discrepancy from positive finite thermal moment systems.
\end{enumerate}
By Theorem~\ref{thm:master}, these are not independent targets; they are charts of the same
source-selection theorem.  The positive research obligation is therefore precise: construct a
completed prime--Gamma estimate that crosses the explicit Abel collar while preserving the Xi
divisor incidence.

\section*{Author contributions and accountability}

Michael Haufschild directed the research and assumes responsibility for the submitted claims.
GPT 5.6 Sol, Fable 5, and Grok 4.5 contributed to proof search, adversarial checking, and
manuscript development.  The human author retains responsibility for correctness, novelty,
exposition, and publication compliance.

\begin{thebibliography}{99}

\bibitem{Bombieri2000}
E.~Bombieri,
\emph{Remarks on Weil's quadratic functional in the theory of prime numbers, I},
Rend. Lincei Mat. Appl. \textbf{11} (2000), 183--233.

\bibitem{Connes1999}
A.~Connes,
\emph{Trace formula in noncommutative geometry and the zeros of the Riemann zeta function},
Selecta Math. (N.S.) \textbf{5} (1999), 29--106.

\bibitem{ConnesConsani2019}
A.~Connes and C.~Consani,
\emph{The scaling Hamiltonian},
arXiv:1910.14368, 2019.

\bibitem{ConnesConsani2021}
A.~Connes and C.~Consani,
\emph{Weil positivity and trace formula, the archimedean place},
arXiv:2006.13771, 2020.

\bibitem{Hausdorff1921}
F.~Hausdorff,
\emph{Summationsmethoden und Momentfolgen. I},
Math. Z. \textbf{9} (1921), 74--109.

\bibitem{Titchmarsh1986}
E.~C.~Titchmarsh,
\emph{The Theory of the Riemann Zeta-Function}, 2nd ed., revised by D.~R.~Heath-Brown,
Oxford University Press, 1986.

\bibitem{Weil1952}
A.~Weil,
\emph{Sur les ``formules explicites'' de la th\'eorie des nombres premiers},
Commun. S\'em. Math. Univ. Lund, volume dedicated to Marcel Riesz (1952), 252--265.

\end{thebibliography}

\end{document}
